% from ega2-errata-addenda-fr.tex
\oldpage[II]{222}
\markboth{A. GROTHENDIECK --- Chap. II}{A. GROTHENDIECK --- Chap. II}
\begin{equation}
\label{ega2:addendum:I.7.3.8.1-fr}
  \tau:f^*(\sh{R}(Y))\to\sh{R}(X).
\tag{7.3.8.1}
\end{equation}

Supposons d'abord $X=\Spec(A)$ et $Y=\Spec(B)$ affines d'anneaux
intègres $A$ et $B$, $f$ correspondant donc à un homomorphisme injectif
$B\to A$, qui se prolonge en un monomorphisme $L\to K$ du corps des
fractions $L$ de $B$ dans le corps des fractions $K$ de $A$.
L'homomorphisme (7.3.8.1) correspond alors à l'homomorphisme canonique
\[
  L\otimes_B A\to K
\]
(I, 1.6.5). Dans le cas général, pour tout couple d'ouverts affines non
vides $U\subset X$, $V\subset Y$ tels que $f(U)\subset V$, on définit
de la façon précédente un homomorphisme $\tau_{U,V}$ et on constate
aussitôt que si $U'\subset U$, $V'\subset V$,
$f(U')\subset V'$, $\tau_{U,V}$ prolonge $\tau_{U',V'}$, d'où notre
assertion. Si $x,y$ sont les points génériques de $X$ et $Y$
respectivement, on a $f(x)=y$,
\[
  \bigl(f^*(\sh{R}(Y))\bigr)_x
  =\sh{O}_y\otimes_{\sh{O}_y}\sh{O}_x
  =\sh{O}_x
\]
(0, 4.3.1), et $\tau_x$ est donc un isomorphisme.

\paragraph{(I, 9.5.2).}
Dans les 3 dernières lignes de la p.~176 et les 5 premières de la
p.~177, remplacer partout $X'$ par $Y'$ et $X''$ par $Y''$.

\paragraph{(I, 10.5.4).}
Ligne 11 du bas de la p.~187, remplacer (10.3.4) par (10.3.5).

\paragraph{(I, 10.6.3).}
Ligne 3 du bas de la p.~189, remplacer $X$ par $\mathfrak X$.

\paragraph{(I, 10.14.2).}
Ligne 5 de la p.~210, remplacer «~$\sh{O}_{\mathfrak X}$-Module cohérent~» par
«~Idéal cohérent de $\sh{O}_{\mathfrak X}$~»; ligne 7, remplacer
«~$\sh{O}_{\mathfrak X}$-Modules cohérents~» par «~Idéaux cohérents de
$\sh{O}_{\mathfrak X}$~».
